
\documentclass[12pt]{article}
\usepackage{amsmath, amsthm, amsfonts, mathtools, xcolor, caption}

% \usepackage{minimalJ}  % decoupled: unused in this file

\numberwithin{equation}{section}

\usepackage[style=numeric,sorting=none,giveninits=true]{biblatex}
\addbibresource{bibliography.bib}

\usepackage{etoolbox}
\usepackage[margin=1in]{geometry}

\DeclareMathAlphabet{\mymathbb}{U}{BOONDOX-ds}{m}{n}

%\usepackage[colorlinks, allcolors=blue, hypertexnames=false]{hyperref}
%\usepackage[noabbrev, nameinlink]{cleveref}

\usepackage{comment}

%%%%%%%%%%%%%%%%%%%%%%%%%%%%%%%%%%%%%%%%%
%       Theorem Environments            %
%%%%%%%%%%%%%%%%%%%%%%%%%%%%%%%%%%%%%%%%%

\newtheorem{theorem}{Theorem}[section]
\newtheorem{lemma}[theorem]{Lemma}
\newtheorem{corollary}[theorem]{Corollary}
\newtheorem{proposition}[theorem]{Proposition}
\theoremstyle{definition}
\newtheorem{remark}{Remark}
\newtheorem{definition}{Definition}
\newtheorem{notation}{Notation}
\newtheorem{example}{Example}
\newtheorem{conjecture}[theorem]{Conjecture}
\newtheorem{openproblem}[theorem]{Open Problem}

%%%%%%%%%%%%%%%%%%%%%%%%%%%%%%%%%%%%%%%%%
%               New Commands            %
%%%%%%%%%%%%%%%%%%%%%%%%%%%%%%%%%%%%%%%%%

\newcommand{\ev}[1]{( #1 )} 
\newcommand{\ew}[1]{\left( #1 \right)} 

\newcommand{\N}{\mathbb{N}}
\newcommand{\Z}{\mathbb{Z}}
\newcommand{\C}{\mathbb{C}}
\newcommand{\Q}{\mathbb{Q}}
\newcommand{\R}{\mathbb{R}}



%%%%%%%%%%%%%%%%%%%%%%%%%%%%%%%%%%%%%%%%%
%               Document                %
%%%%%%%%%%%%%%%%%%%%%%%%%%%%%%%%%%%%%%%%%

\let\underbrace\LaTeXunderbrace
\let\overbrace\LaTeXoverbrace


\begin{document}

\section*{Challenge 10: The Unfriendly Partition Conjecture}
\refstepcounter{section}

(Cowan and Emerson, see for example \cite{aharoni1990unfriendly})


\begin{definition}[Partial partition]
A {\em partial partition} of the vertices of a graph $G$ is a partial function $V(G) \to \{0,1\}$. We call it {\em finite} if its domain is finite. We call it a {\em partition} if it is a total function.
\end{definition}

\begin{definition}[Unfriendly]
Let $v$ be a vertex of a graph $G$. A partial partition $f \colon V(G) \to \{0,1\}$ is {\em unfriendly at $v$} if $v$ has at least as many neighbours mapped by $f$ to the opposite partition class as it has neighbours mapped to the same partition class. Formally this means
$$|N(v) \cap f^{-1}[1 - f(v)]| \ge |N(v) \cap f^{-1}[f(v)]|.$$
A partition is {\em unfriendly} if it is unfriendly at every vertex.
\end{definition}

\begin{conjecture}[Unfriendly partition Conjecture]
Every countable graph has an unfriendly partition.
\end{conjecture}

In order to pose a scaled version of this conjecture, we introduce the following game:

\begin{definition}[The $r$-partitioning game.]
Let $r$ be an ordinal number and let $G$ be a graph. The {\em $r$-partitioning game on $G$} is played between two players, Challenger and Partitioner, as follows. A position in the game is a pair $(f, \alpha)$ consisting of a partial partition $f$ of $G$ and an ordinal number $\alpha$. The initial position is $(\emptyset, r)$. In each turn, beginning in position $(f, \alpha)$, Challenger must name a vertex $v$ of $G$ and an ordinal $\beta < \alpha$ and Partitioner must then name a finite partial partition $g$ extending $f$ which is defined at $v$ and unfriendly at $v$. After this, the new position is $(g, \beta)$. If at any point some player is unable to move, the other player wins.
\end{definition}

Note that since there are no infinite decreasing sequences of ordinals this game must end after finitely many turns. 

The scaled conjecture is as follows:

\begin{conjecture}[Scaled Unfriendly Partition Conjecture]
Let $r$ be an ordinal. For any countable graph $G$, Partitioner has a winning strategy in the $r$-partitioning game on $G$.
\end{conjecture}

%%%%%%%%%%%%%%%%%%%%%%%%%%%%%%%%%%%%%%%%%%%%%%%%%%
\printbibliography
%%%%%%%%%%%%%%%%%%%%%%%%%%%%%%%%%%%%%%%%%%%%%%%%%%
\end{document}
